% Canonical Paper IV section.
The projective quotient in Section~\ref{sec:condensation-continuation} is one
instance of a more general construction.  A past is relevant only through the
effects it can still have on legal futures.  The definition must therefore
name the interface, the future class, the observation, and the treatment of
partial domains.

\subsection{Typed interfaces and totalized partial action}

Fix a typed interface $C$.  Let $P_C$ be the set of admissible pasts reaching
that interface and let $M_C$ be a monoid of future continuations.  Composition
is written as a right action $p\cdot\eta$ and may initially be partial because
ordinary arithmetic, typing, or resource rules can make a continuation
illegal.

If the action is partial, introduce an absorbing state $\bot$ and put
\[
  P_C^\bot=P_C\sqcup\{\bot\}.
\]
If the action is already total, use the convention $P_C^\bot=P_C$ and adjoin
nothing.  In the partial case, totalize the action by declaring an illegal
continuation to have value $\bot$ and by setting
$\bot\cdot\eta=\bot$.  Let
\[
  O_C:P_C^\bot\longrightarrow Y_C^\bot
\]
be the current observation, with $O_C(\bot)=\bot$ when the absorbing state
is present.  Depending on the problem,
$O_C$ may record an endpoint, a projective state, a Boolean answer, a tensor, a
trace, or a distribution.

\begin{definition}[Contextual continuation equivalence]
\label{def:contextual-equivalence}
For $p,p'\in P_C^\bot$, define
\begin{equation}
  p\equiv_C p'
  \quad\Longleftrightarrow\quad
  O_C(p\cdot\eta)=O_C(p'\cdot\eta)
  \quad\text{for every }\eta\in M_C.
  \label{eq:contextual-equivalence}
\end{equation}
The quotient
\[
  \cR_C=P_C^\bot/{\equiv_C}
\]
is the contextual continuation-residual space at $C$.
\end{definition}

The absorbing outcome is essential for partial arithmetic.  Two pasts are not
identified merely because all common legal continuations agree if one of them
admits a future that the other does not.

\begin{proposition}[Equivalence, refinement, and observation monotonicity]
\label{prop:residual-basic-properties}
The relation $\equiv_C$ is an equivalence relation.  Enlarging the continuation
class can only refine it.  Replacing $O_C$ by a coarser observation
$f\circ O_C$ can only coarsen it.
\end{proposition}

\begin{proof}
Reflexivity, symmetry, and transitivity hold pointwise for every continuation.
A larger family supplies more equality tests.  Applying a function to equal
observations preserves equality and can identify previously distinct values.
\end{proof}

Thus a residual is never absolute.  It is relative to a declared experimental
interface.

\subsection{Residuals form the canonical continuation machine}

\begin{proposition}[Right congruence]
\label{prop:residual-right-congruence}
If $p\equiv_Cp'$, then for every $\zeta\in M_C$,
\[
  p\cdot\zeta\equiv_Cp'\cdot\zeta.
\]
Consequently the formula
\[
  [p]\cdot\zeta=[p\cdot\zeta]
\]
defines a right action of $M_C$ on $\cR_C$.
\end{proposition}

\begin{proof}
For every further continuation $\eta$,
\[
  O_C((p\cdot\zeta)\cdot\eta)
  =O_C(p\cdot(\zeta\eta))
  =O_C(p'\cdot(\zeta\eta))
  =O_C((p'\cdot\zeta)\cdot\eta).
\]
\end{proof}

The identity continuation shows that
\[
  \overline O_C([p])=O_C(p)
\]
is well defined.  Hence $(\cR_C,M_C,\overline O_C)$ is itself an exact
deterministic continuation machine.

\begin{definition}[Exact deterministic online realization]
An exact deterministic online realization of the experiment at $C$ consists
of a state set $S$, a right action of $M_C$ on $S$, an encoding
$e:P_C^\bot\to S$, and an output map $o:S\to Y_C^\bot$ such that
\begin{equation}
  e(p\cdot\eta)=e(p)\cdot\eta,
  \qquad
  o(e(p))=O_C(p).
  \label{eq:exact-machine-equivariance}
\end{equation}
The first identity includes the totalized invalid state when necessary.
\end{definition}

The equivariance requirement prevents a representation from being called an
online state merely because a separate external procedure knows the entire
past and recomputes transitions from it.

\begin{theorem}[Canonical minimal exact state]
\label{thm:minimal-exact-residual}
For every exact deterministic online realization there is a unique surjective
$M_C$-equivariant map
\[
  \phi:e(P_C^\bot)\longrightarrow\cR_C
  \qquad\text{with}\qquad
  \phi(e(p))=[p].
\]
It also preserves outputs.  Therefore $\cR_C$ is the canonical minimal
reachable exact state space, up to isomorphism.
\end{theorem}

\begin{proof}
If $e(p)=e(p')$, then for every $\eta$ equivariance and output correctness give
\[
 O_C(p\cdot\eta)
 =o(e(p)\cdot\eta)
 =o(e(p')\cdot\eta)
 =O_C(p'\cdot\eta).
\]
Hence $p\equiv_Cp'$, so $\phi$ is well defined.  It is surjective by definition
of $\cR_C$.  Equivariance follows from
Proposition~\ref{prop:residual-right-congruence}, and output preservation follows
from the identity continuation.  Uniqueness is forced on every reachable state
$e(p)$.
\end{proof}

This is the finite-state principle behind the classical theory of residual
languages and automata~\cite{Nerode1958}, but the present formulation permits
arbitrary typed observations and partial arithmetic domains.

\subsection{Finite-state information lower bound}

\begin{corollary}[Fixed-width exact-state bound]
\label{cor:residual-bit-bound}
Assume $\cR_C$ is finite.  If an exact reachable state is stored in a fixed
$b$-bit register, then
\begin{equation}
  b\ge \left\lceil\log_2|\cR_C|\right\rceil.
  \label{eq:residual-bit-bound}
\end{equation}
The same conclusion holds for a prefix-free variable-length encoding whose
maximum codeword length is $b$.
\end{corollary}

\begin{proof}
Theorem~\ref{thm:minimal-exact-residual} requires at least $|\cR_C|$ reachable
code states.  A fixed $b$-bit register has at most $2^b$ states.  For a
prefix-free code of maximum length $b$, Kraft's inequality also gives at most
$2^b$ codewords.
\end{proof}

The theorem is a state-selection lower bound.  It does not determine the size
of a transition table, the cost of computing the encoding, the time needed to
apply a continuation, or the materialized size of a structured value.

\begin{corollary}[Fixed-register capacity--time bound]
\label{cor:capacity-time-bound}
Consider declared cuts $C_0,\ldots,C_T$.  Suppose that at cut $C_t$ a complete
exact state for every admissible past must fit in a fixed-width register of
capacity $b_t$ bits, and charge one full register snapshot at each cut.  Then
\begin{equation}
  \sum_{t=0}^{T} b_t
  \ge
  \sum_{t=0}^{T}
  \left\lceil\log_2|\cR_{C_t}|\right\rceil.
  \label{eq:capacity-time-bound}
\end{equation}
If $|\cR_{C_t}|\ge c t^d$ eventually, the right side is
$\Omega(T\log T)$; if $|\cR_{C_t}|\ge c a^t$ with $a>1$, it is
$\Omega(T^2)$.
\end{corollary}

\begin{proof}
Apply Corollary~\ref{cor:residual-bit-bound} at each cut and sum.  The two
asymptotic statements follow from summing $\log t$ and $t$, respectively.
\end{proof}

Equation~\eqref{eq:capacity-time-bound} is deliberately called a
capacity--time bound.  It becomes an actual memory--time bound for a run only
when the operational model charges the complete register as live memory at
each snapshot.  Variable-size structured states, compression, communication,
and recomputation require the model of Section~\ref{sec:operational-resources}.

\subsection{Three projective residual experiments}

The quotient tower supplies three different experiments.

\begin{enumerate}
  \item If pasts are operators $g\in G$, futures act on the left, and the
  observation is the bivaluation $gH$, then
  $\cR\cong G/H$ by Corollary~\ref{cor:left-continuation-exact}.

  \item If the observation remembers only $gq_0^+$ and futures continue to act
  on the left, then the residual is $G/B_+$.  The kernel direction has become
  observationally irrelevant.

  \item If arbitrary reversible futures act on the right, $G/H$ does not
  support the action unless the futures lie in $N_G(H)$.  A full operator lift,
  or another state carrying the missing frame information, is then required.
\end{enumerate}

Thus projective process residue can coincide with a continuation residual, map
to one, or be necessary to refine one.  Which case occurs is a theorem about
the chosen action and observation, not a property of the static fiber alone.

\subsection{Randomized and approximate boundaries}

Definition~\ref{def:contextual-equivalence} uses exact deterministic equality.
A randomized algorithm may permit equal output distributions rather than equal
outputs; an approximate algorithm may identify observations within a metric or
loss tolerance; a statistical model may retain only a sufficient statistic for
a declared distribution family.  Each change alters the residual relation and
its lower bound.  No such approximation theorem is inferred from
Corollary~\ref{cor:residual-bit-bound}; approximate residual geometry remains an
open programme.
